\chapter{フォルダ構成と配置場所}

\section{配置先のディレクトリ}

\cmd{git clone} を実行すると、ホームディレクトリ直下に
リポジトリフォルダが作成されます。
課題フォルダはリポジトリ直下の \cmd{2026年度\_B4課題引継ぎ/} にあります。

\begin{wslbox}[フォルダ構成]
\begin{lstlisting}[style=bash]
/home/username/              <- ホームディレクトリ（~ と同じ）
└── NMR-lab-kadai/           <- git clone で作成される
    ├── 2026年度_B4課題引継ぎ/  <- VS Code で開くフォルダ
    │   ├── .vscode/
    │   │   └── settings.json  <- カーネル自動設定
    │   ├── pyproject.toml
    │   ├── uv.lock
    │   ├── No.1_FID/
    │   │   ├── makesignal.ipynb
    │   │   └── heikatsu.ipynb
    │   ├── No.2_FFT1D/
    │   │   ├── FFT1D.ipynb
    │   │   ├── kukei.ipynb
    │   │   └── sft.ipynb
    │   ├── No.3_FFT2D/
    │   │   └── FFT2D.ipynb
    │   ├── No.4_DFT_exercise/
    │   │   └── DFT_template.ipynb
    │   ├── No.5_CompressedSensing/
    │   │   ├── cs_tv.ipynb
    │   │   └── cs_wavelet.ipynb
    │   ├── No.6_Denoising/
    │   │   ├── add_noise.ipynb
    │   │   └── denoise_bm3d.ipynb
    │   ├── No.7_Metrics/
    │   │   └── psnr_ssim.ipynb
    │   └── No.8_CT/
    │       └── ct_pipeline.ipynb
    └── uv-manual/             <- セットアップマニュアル
\end{lstlisting}
\end{wslbox}

\begin{notebox}[username について]
\cmd{username} の部分は自分の WSL ユーザー名です。
次のコマンドで確認できます。

\begin{wslbox}[ユーザー名の確認]
\begin{lstlisting}[style=bash]
$ whoami
\end{lstlisting}
\end{wslbox}
\end{notebox}

\section{重要なファイルの説明}

\begin{center}
\begin{tabularx}{\linewidth}{lX}
  \toprule
  ファイル & 役割 \\
  \midrule
  \cmd{pyproject.toml} &
    必要なライブラリとバージョン条件が書かれた設定ファイル。
    \cmd{uv sync} はこれを読んでインストール内容を決める \\
  \cmd{uv.lock} &
    実際にインストールするバージョンを厳密に記録したファイル。
    全員が同じ環境を再現するために使う。
    \textbf{編集しないこと} \\
  \cmd{.venv/} &
    \cmd{uv sync} を実行すると自動生成される仮想環境フォルダ。
    git には含まれないので各自が \cmd{uv sync} で作成する \\
  \cmd{.vscode/settings.json} &
    VS Code が自動で \cmd{.venv} の Python を選ぶための設定ファイル。
    カーネル選択の手間を省く \\
  \cmd{*.ipynb} &
    JupyterLab で開く課題ノートブック \\
  \cmd{*.py} &
    対応する Python スクリプト（ノートブックと同じ処理） \\
  \bottomrule
\end{tabularx}
\end{center}

\begin{notebox}[.venv/ は git 管理外]
\cmd{.venv/} フォルダはサイズが数百 MB あり、
\cmd{uv sync} で誰でも再生成できるため git には含まれていません。
\cmd{git clone} 後に各自で \cmd{uv sync} を実行してください。
\end{notebox}
